\documentclass[UTF8,a4paper,11pt]{ctexart}

\usepackage[a4paper,margin=2cm]{geometry}
\usepackage{array}
\usepackage{booktabs}
\usepackage{enumitem}
\usepackage{fancyhdr}
\usepackage{hyperref}
\usepackage{listings}
\usepackage[table]{xcolor}
\usepackage{longtable}
\usepackage{tabularx}
\usepackage{tikz}
\usepackage[most]{tcolorbox}
\usepackage{titlesec}
\usepackage{etoolbox}
\usetikzlibrary{arrows.meta,positioning,shapes.geometric}

\hypersetup{
  colorlinks=true,
  linkcolor=Summary,
  urlcolor=blue!60!black
}

\definecolor{Accent}{HTML}{1F5C4A}
\definecolor{AccentLight}{HTML}{E7F1EE}
\definecolor{Summary}{HTML}{355C8C}
\definecolor{SummaryLight}{HTML}{EAF0F8}
\definecolor{Warn}{HTML}{A34D32}
\definecolor{WarnLight}{HTML}{FCEEE9}
\definecolor{CodeBg}{HTML}{F6F8F9}
\definecolor{GrayLine}{HTML}{D8DEE1}
\definecolor{CoverSand}{HTML}{F4E8D8}

\pagestyle{fancy}
\fancyhf{}
\fancyhead[L]{Codex CLI 讲义摘要}
\fancyhead[R]{速读版}
\fancyfoot[C]{\thepage}
\setlength{\headheight}{15pt}

\titleformat{\section}[display]
  {\LARGE\bfseries\color{Accent}}
  {\color{Summary}\rule{\textwidth}{1pt}\\[-0.25em]\large\thesection}
  {0.3em}
  {\filright}
\titlespacing*{\section}{0pt}{0pt}{0.9em}
\titleformat{\subsection}{\large\bfseries\color{Accent}}{\thesubsection}{0.6em}{}
\setlist[itemize]{leftmargin=1.6em,noitemsep,topsep=0.25em}
\setlist[enumerate]{leftmargin=1.6em,noitemsep,topsep=0.25em}
\setlength{\parindent}{2em}
\setlength{\parskip}{0.12em}

\lstdefinestyle{cli}{
  backgroundcolor=\color{CodeBg},
  basicstyle=\ttfamily\small,
  frame=single,
  rulecolor=\color{GrayLine},
  columns=fullflexible,
  breaklines=true,
  breakautoindent=false,
  breakindent=0pt,
  showstringspaces=false
}

\lstdefinestyle{prompt}{
  style=cli,
  linewidth=0.9\linewidth,
  xleftmargin=0.05\linewidth
}

\newtcolorbox{tipbox}{
  colback=SummaryLight,
  colframe=Summary,
  boxrule=0.7pt,
  arc=2mm,
  title=记住这一点,
  fonttitle=\bfseries
}

\newtcolorbox{warnbox}{
  colback=WarnLight,
  colframe=Warn,
  boxrule=0.7pt,
  arc=2mm,
  title=别踩这个坑,
  fonttitle=\bfseries
}

\newtcolorbox{sumbox}{
  colback=AccentLight,
  colframe=Accent,
  boxrule=0.7pt,
  arc=2mm,
  title=速记,
  fonttitle=\bfseries
}

\pretocmd{\section}{\clearpage}{}{}

\begin{document}

\thispagestyle{empty}
\begin{tikzpicture}[remember picture,overlay]
  \fill[Accent!10] ([xshift=-1.3cm,yshift=1.0cm]current page.north east) circle (3.1cm);
  \fill[Summary!10] ([xshift=0.7cm,yshift=-0.5cm]current page.north west) circle (2.0cm);
  \fill[CoverSand] ([xshift=-1.0cm,yshift=-0.8cm]current page.south east) circle (2.5cm);
\end{tikzpicture}

\vspace*{0.6cm}

\begin{tcolorbox}[
  enhanced,
  colback=white,
  colframe=Accent,
  boxrule=1pt,
  arc=3mm,
  left=7mm,
  right=7mm,
  top=7mm,
  bottom=7mm
]
{\Large\bfseries\color{Accent}Codex CLI 讲义摘要}\\[0.35em]
{\large\color{black!72}适合培训、复习、分享的速读版}\\[0.7em]
\small
\tcbox[colback=AccentLight,colframe=Accent,arc=2mm,boxrule=0.7pt,left=2mm,right=2mm,top=0.8mm,bottom=0.8mm]{第一次上手}
\hspace{0.5em}
\tcbox[colback=SummaryLight,colframe=Summary,arc=2mm,boxrule=0.7pt,left=2mm,right=2mm,top=0.8mm,bottom=0.8mm]{权限与审批}
\hspace{0.5em}
\tcbox[colback=WarnLight,colframe=Warn,arc=2mm,boxrule=0.7pt,left=2mm,right=2mm,top=0.8mm,bottom=0.8mm]{提示词速记}

\vspace{0.7em}

这份短版保留了最核心的起步路径：安装、提问、权限判断、常用命令、稳定工作流、排错顺序和最后的速记卡。
\end{tcolorbox}

\begin{tcolorbox}[
  colback=CodeBg,
  colframe=GrayLine,
  boxrule=0.6pt,
  arc=2mm,
  left=5mm,
  right=5mm,
  top=3mm,
  bottom=3mm
]
\small
请以你本地 \texttt{codex --help} 和 OpenAI 官方 Codex 文档为准。
\end{tcolorbox}

\section{一页理解}

先记住一句最实用的话：\textbf{Codex CLI 是一个在终端里工作的 AI 编码代理。}

它最适合做四类事：
\begin{itemize}
  \item 看懂陌生项目；
  \item 定位功能和报错；
  \item 做小范围改动；
  \item 做代码审查、总结、给出验证建议。
\end{itemize}

\begin{center}
\begin{tikzpicture}[
  box/.style={draw=Accent, rounded corners, fill=AccentLight, minimum width=2.7cm, minimum height=0.95cm, align=center},
  line/.style={-Latex, thick, draw=Accent},
  node distance=8mm and 7mm
]
\node[box] (a) {进入项目目录};
\node[box, right=of a] (b) {说清楚任务};
\node[box, right=of b] (c) {让它读上下文};
\node[box, below=of b] (d) {看解释或改动};
\node[box, left=of d] (e) {继续追问};
\node[box, right=of d] (f) {复核与验证};
\draw[line] (a) -- (b);
\draw[line] (b) -- (c);
\draw[line] (c) -- (d);
\draw[line] (d) -- (e);
\draw[line] (d) -- (f);
\end{tikzpicture}
\end{center}

\begin{sumbox}
第一次用 Codex CLI，不求全会。先掌握一条线：进目录、提具体任务、继续追问、最后验证。
\end{sumbox}

\section{第一次上手}

\subsection{安装与登录}

\begin{lstlisting}[style=cli]
npm i -g @openai/codex
codex --help
codex
\end{lstlisting}

如果你想显式登录或检查状态：

\begin{lstlisting}[style=cli]
codex login
codex login status
\end{lstlisting}

如果你走 API 密钥路线：

\begin{lstlisting}[style=cli]
printenv OPENAI_API_KEY | codex login --with-api-key
\end{lstlisting}

\begin{tipbox}
这两个变量名别混：上面这条命令是在“登录”时从 \texttt{OPENAI\_API\_KEY} 读入密钥；到了 \texttt{codex exec} 和 CI 场景，官方更常用的是 \texttt{CODEX\_API\_KEY}。
\end{tipbox}

\subsection{Linux 环境下先确认这四件事}

下面默认你用的是 Linux 终端。无论你是在本地、SSH 到远端，还是在容器里，第一次开始前都建议先跑：

\begin{lstlisting}[style=cli]
pwd
ls
codex --help
codex login status
\end{lstlisting}

这四步分别对应四个关键问题：我现在在哪、这是不是目标项目、CLI 能不能用、认证是不是已经就绪。

\subsection{macOS、Windows、WSL 一句话提醒}

\begin{itemize}
  \item \textbf{macOS}：整体和 Linux 很像，先确认 \texttt{codex} 在 PATH 里。
  \item \textbf{PowerShell}：环境变量写法和 Bash 不一样，别直接照抄。
  \item \textbf{WSL}：把它当独立 Linux；如果项目主要在 WSL 里，就尽量在 WSL 里完成登录、Git 和测试。
\end{itemize}

\begin{warnbox}
最容易乱的是 Windows 和 WSL 混着用。项目主要在哪边，就尽量整轮任务都在哪边完成。
\end{warnbox}

\subsection{第一次提问}

\begin{lstlisting}[style=prompt]
请先看一下当前项目结构，告诉我入口文件、核心目录和推荐阅读顺序，并请引用具体文件路径。
\end{lstlisting}

\begin{warnbox}
不要一上来就说“帮我改好”或“帮我优化一下”。范围太宽，结果通常会变空。
\end{warnbox}

\section{最常用命令}

把第一句任务说清楚以后，下一步就是认入口：什么时候直接进会话，什么时候跑单次命令，什么时候留在会话里用斜杠命令。先别把这一页当命令大全来背，先把“眼前这个场景该走哪条路”搞明白。

\rowcolors{2}{white}{AccentLight}
\begin{longtable}{>{\ttfamily\raggedright\arraybackslash}p{0.24\textwidth} >{\raggedright\arraybackslash}p{0.22\textwidth} >{\raggedright\arraybackslash}p{0.40\textwidth}}
\toprule
命令 & 作用 & 什么时候用 \\
\midrule
\endhead
codex & 交互式会话 & 本地日常开发最常用 \\
codex exec "..." & 非交互执行 & 脚本、CI（持续集成）、一次性任务 \\
codex review --uncommitted & 审查当前改动 & 提交前快速找风险 \\
codex resume --last & 接着上次会话干 & 不想重复讲背景时 \\
codex apply <TASK\_ID> & 应用某次任务 diff & 结果要落到工作树时 \\
/new & 开一轮新会话 & 不退出终端，直接重开一轮 \\
/clear & 清屏并新开会话 & 想清掉当前屏幕内容时 \\
/resume & 恢复旧会话 & 想切回之前保存的会话 \\
/fork & 从旧会话分叉 & 想保留上下文，但开一条新线 \\
/status & 查看会话状态 & 看模型、权限、可写目录 \\
/permissions & 调整权限模式 & 切换更保守或更自动 \\
/review & 会话内做代码审查 & 先查风险再看 diff \\
/diff & 查看改动 & 改完先审一眼 \\
/plan & 先出计划 & 大任务先计划再动手 \\
/init & 生成 \texttt{AGENTS.md} 脚手架 & 想沉淀仓库规则时 \\
\bottomrule
\end{longtable}

\begin{tipbox}
如果你只背三个入口，先背 \texttt{codex}、\texttt{codex exec}、\texttt{codex review}。
\end{tipbox}

\subsection{现在该用哪个入口？}

\begin{itemize}
  \item 想连续对话、边看边追问：\texttt{codex}
  \item 想脚本化、跑 CI、做一次性任务：\texttt{codex exec}
  \item 已经在会话里，想原地重开：\texttt{/new} 或 \texttt{/clear}
  \item 想回到旧会话：\texttt{/resume} 或 \texttt{codex resume --last}
  \item 想保留旧思路，再开一条新线：\texttt{/fork}
\end{itemize}

\subsection{什么时候开 \texttt{--search}}

\begin{itemize}
  \item 只看本地代码：项目结构、调用链、当前分支改动、仓库里的测试命令。
  \item 开 \texttt{--search}：官方文档、依赖版本差异、外部 API、最新错误说明。
\end{itemize}

\begin{tipbox}
先问本地事实，再补外部资料。并要求把“仓库事实”和“外部来源”分开写。
\end{tipbox}

\section{权限与审批}

\subsection{三个沙箱级别}

\begin{itemize}
  \item \texttt{read-only}：只读，不改。
  \item \texttt{workspace-write}：能改当前工作区，最适合本地开发。
  \item \texttt{danger-full-access}：权限很大，只在你非常清楚风险时使用。
\end{itemize}

\subsection{三个审批策略}

\begin{itemize}
  \item \texttt{untrusted}：安全读操作自动跑，不信任命令先问你。
  \item \texttt{on-request}：由代理判断什么时候该申请批准。
  \item \texttt{never}：不弹审批，常用于自动化，但前提是你自己已经控好风险。
\end{itemize}

\begin{lstlisting}[style=cli]
codex --full-auto
\end{lstlisting}

它相当于一个比较省事的默认组合：

\begin{lstlisting}[style=cli]
--sandbox workspace-write --ask-for-approval on-request
\end{lstlisting}

\begin{sumbox}
本地开发优先从 \texttt{workspace-write} 起步。先别急着全放开。
\end{sumbox}

\section{提示词速记}

\subsection{一条好提示词的四要素}

\begin{enumerate}
  \item 目标：你到底要它做什么。
  \item 范围：只围绕哪些文件、功能或问题。
  \item 限制：不要做什么。
  \item 输出格式：你希望它最后怎么汇报。
\end{enumerate}

\subsection{直接可用模板}

\textbf{看项目：}
\begin{lstlisting}[style=prompt]
请按新手视角介绍当前项目，重点说入口、核心目录、主流程和推荐阅读顺序。
\end{lstlisting}

\textbf{定位功能：}
\begin{lstlisting}[style=prompt]
请帮我定位登录相关代码，告诉我入口、涉及文件和大致调用链。
\end{lstlisting}

\textbf{处理报错：}
\begin{lstlisting}[style=prompt]
这是报错信息。请先解释最可能原因，再给排查顺序，最后给最小修改方案。
\end{lstlisting}

\textbf{小改动：}
\begin{lstlisting}[style=prompt]
请只修改提示文案，不要改业务逻辑。
修改后列出改动文件和验证方法。
\end{lstlisting}

\section{稳定工作流}

这一节是把前面那些零散动作串成一条主线：先理解、再定位、再小改、最后复核。后面的“会话控制”和“自动化与 CI”更像按需补充，不必第一次就全背下来。

\subsection{本地开发推荐顺序}

\begin{enumerate}
  \item 先解释项目或问题；
  \item 再定位文件和调用链；
  \item 再做小范围修改；
  \item 改完跑 \texttt{/review}；
  \item 再看 \texttt{/diff}；
  \item 最后自己验证。
\end{enumerate}

\subsection{改坏了先怎么救}

\begin{lstlisting}[style=cli]
git status
git diff
git switch -c backup/codex-check
git stash push -u -m "before-rollback"
\end{lstlisting}

先看现场，再保留现场，最后才决定撤哪一部分。能只撤一个文件，就别一上来整批清空。

\subsection{三轮追问示例}

\begin{enumerate}
  \item 第一轮拿地图：先问入口文件、相关目录、阅读顺序。
  \item 第二轮抓主线：只盯一个功能，让它列关键文件和调用链。
  \item 第三轮要交付：只改一个小点，并要求按“改动文件 / 验证 / 风险”总结。
\end{enumerate}

\subsection{任务完成验收卡}

一轮任务结束前，至少自己确认四件事：
\begin{itemize}
  \item 改动范围有没有控住；
  \item 实际跑了哪些验证；
  \item 还有什么剩余风险；
  \item 现在适不适合提交。
\end{itemize}

如果你想让 Codex 按这个格式汇报，可以直接说：

\begin{lstlisting}[style=prompt]
请按“改动文件 / 验证情况 / 剩余风险 / 是否建议现在提交”四项总结。
\end{lstlisting}

\subsection{为什么要用 \texttt{AGENTS.md}}

\texttt{AGENTS.md} 适合沉淀这些长期规则：
\begin{itemize}
  \item 仓库结构和重点目录；
  \item 构建、测试、静态检查命令；
  \item 哪些文件不要动；
  \item 交付时怎么总结。
\end{itemize}

\begin{lstlisting}[style=cli]
/init
\end{lstlisting}

\begin{tipbox}
如果你在三次任务里都重复说过同一句话，那它大概率就该进 \texttt{AGENTS.md} 了。
\end{tipbox}

\subsection{个人默认习惯写进 \texttt{config.toml}}

\texttt{AGENTS.md} 管仓库规则；\texttt{config.toml} 管你自己的默认习惯，比如常用模型、审批和沙箱组合。

\begin{lstlisting}[style=cli]
codex -p safe
codex -p dev
codex --oss --local-provider ollama -p local
\end{lstlisting}

先想清楚你最常用的 2 到 3 种模式，再去固化。

\section{会话控制}

\subsection{上下文太长怎么办？}

优先做四件事：
\begin{enumerate}
  \item 用 \texttt{/compact} 压缩长会话；
  \item 用 \texttt{/mention} 点名关键文件；
  \item 把任务缩到一个功能或一个报错；
  \item 真想重开一轮，就先用 \texttt{/new}；想清屏就用 \texttt{/clear}。
\end{enumerate}

\subsection{继续旧会话和新开会话怎么选？}

\begin{itemize}
  \item 想原地开启新会话：\texttt{/new}
  \item 想清屏后再开：\texttt{/clear}
  \item 想切回旧会话：\texttt{/resume}
  \item 想基于旧会话分一条新线：\texttt{/fork}
  \item 想从命令行里直接接着最近一次聊：\texttt{codex resume --last}
\end{itemize}

\subsection{能在多个会话间切换吗？}

可以，但更像“恢复某个已保存会话”，不是像浏览器标签页那样一直挂着来回切换。日常最常用的是 \texttt{/resume}、\texttt{/fork}，以及命令行里的 \texttt{codex resume}。

\subsection{对照图}

\begin{center}
\begin{tcbraster}[raster columns=2, raster equal height=rows, raster column skip=6pt, raster row skip=6pt]
\begin{tcolorbox}[colback=AccentLight, colframe=Accent, title=\texttt{/new}]
原地开新会话。\\
\textbf{适合：}不退出终端，直接重开。
\end{tcolorbox}
\begin{tcolorbox}[colback=SummaryLight, colframe=Summary, title=\texttt{/clear}]
清屏后再开新会话。\\
\textbf{适合：}屏幕太乱，想顺手清掉。
\end{tcolorbox}
\begin{tcolorbox}[colback=AccentLight, colframe=Accent, title=\texttt{/resume}]
恢复旧会话。\\
\textbf{适合：}切回上次的上下文。
\end{tcolorbox}
\begin{tcolorbox}[colback=SummaryLight, colframe=Summary, title=\texttt{/fork}]
从旧会话分叉。\\
\textbf{适合：}保留旧讨论，再开新线。
\end{tcolorbox}
\end{tcbraster}
\end{center}

\section{自动化与 CI}

前面讲的，都是你人在终端前盯着会话用的情况。下面切到另一类场景：不连续对话，而是把一次性任务塞进脚本或 CI。

\subsection{最小示例}

\begin{lstlisting}[style=cli]
export CODEX_API_KEY=your_api_key

codex exec --json \
  -o codex-summary.txt \
  "审查当前仓库，并总结最主要的 3 个风险点"
\end{lstlisting}

\subsection{这三个参数分别干什么}

\begin{itemize}
  \item \texttt{CODEX\_API\_KEY}：给非交互执行做认证。
  \item \texttt{--json}：让输出更适合脚本读取。
  \item \texttt{-o}：把最后一条摘要单独写到文件里。
\end{itemize}

\subsection{脚本闭环怎么搭}

\begin{lstlisting}[style=cli]
codex exec --json \
  --output-schema review-schema.json \
  -o codex-summary.txt \
  "审查当前改动" > codex-events.jsonl
\end{lstlisting}

给人看：\texttt{codex-summary.txt}。给脚本读：\texttt{codex-events.jsonl}。要稳定结构：\texttt{--output-schema}。

\begin{warnbox}
自动化更适合做总结、分类、只读检查。大范围自动改代码，先在本地跑稳再说。
\end{warnbox}

\section{排错与避坑}

\subsection{真实排错顺序}

当你遇到报错时，最稳的顺序通常是：
\begin{enumerate}
  \item 先让它解释错误类型；
  \item 再让它定位最可能的文件；
  \item 再让它给最小修复方案；
  \item 最后让它说明验证方法。
\end{enumerate}

\subsection{最容易踩的坑}

\begin{itemize}
  \item 没确认当前目录就开始操作。
  \item 任务太宽，结果很空。
  \item 一上来就给太大权限。
  \item 改完不看 \texttt{/diff}，也不做验证。
  \item 在 CI 里直接照搬交互式习惯。
\end{itemize}

\subsection{高风险边界}

下面这些事，别一上来就让 Codex 直接放手执行：
\begin{itemize}
  \item 删除、覆盖、批量改名、批量格式化；
  \item 破坏性 Git 操作；
  \item 数据库迁移、数据回填、修复脚本；
  \item 密钥、认证缓存、生产配置；
  \item 会触发外部副作用的命令。
\end{itemize}

更稳的做法是先说：先不要执行，先列风险、影响范围、回滚方法。

\begin{sumbox}
新手最大的问题不是不会敲命令，而是没有形成节奏：确认目录、收紧范围、审改动、做验证。
\end{sumbox}

\section{FAQ}

\subsection{上下文窗口超了怎么办？}

不要继续硬塞更多内容，优先这样做：
\begin{enumerate}
  \item 缩小任务范围，只盯一个功能或一个报错。
  \item 用 \texttt{/compact} 压缩长会话。
  \item 只点名真正关键的文件，不要一次拉太多目录。
  \item 必要时开新会话，并先用几句话总结上一轮结论。
\end{enumerate}

\subsection{它回答得很空怎么办？}

大多数时候是任务太宽。把问题改成“哪几个文件、哪个功能、最后要什么格式”，通常会立刻好很多。

\subsection{最少要做哪些复核？}

如果你只做最低限度检查，至少做三件事：看 \texttt{/diff}、跑相关测试或静态检查、自己验证关键路径。时间允许的话，再补一次 \texttt{/review} 会更稳。

\section{高手常用招式}

到这里，入门主线其实已经够你正常开工了。下面这部分更像“用熟以后再加上的控场技巧”，适合按需查，不用第一次就全部照做。

\subsection{这 8 招最值得记}

\begin{enumerate}
  \item 先比方案，再决定改不改。
  \item 大仓库先收文件范围，不要一口气全讲。
  \item 用 \texttt{/fork} 保留旧结论，再试另一条线。
  \item 用 \texttt{/model} 做“先快后深”的分工。
  \item 先定义成功标准，再开始修改。
  \item 把代码审查做成闭环，不是只跑一次。
  \item 固定交付模板，别让结果每次都飘。
  \item \texttt{codex exec} 优先做总结、分类、只读检查。
\end{enumerate}

\subsection{直接可抄的进阶提示词}

\textbf{方案评审：}
\begin{lstlisting}[style=prompt]
先不要改代码。请给我 2 到 3 种实现方案，分别说明改动范围、优点、风险，并给推荐。
\end{lstlisting}

\textbf{大仓库控上下文：}
\begin{lstlisting}[style=prompt]
先不要全仓库泛讲。请先找出和“登录失败提示”最相关的 5 个文件，只输出路径，并告诉我下一轮先追哪两个。
\end{lstlisting}

\textbf{成功标准前置：}
\begin{lstlisting}[style=prompt]
先不要修改。请先复述目标、列出不做的事、给出完成标准和验证方法。等我确认后再开始。
\end{lstlisting}

\textbf{结构化交付：}
\begin{lstlisting}[style=prompt]
请固定按“结论 / 改动文件 / 验证情况 / 剩余风险 / 下一步建议”汇报。
\end{lstlisting}

\subsection{最常见的高阶误用}

\begin{itemize}
  \item 明明该先比方案，却一上来就让它直接改。
  \item 明明仓库很大，却还在不断塞更多上下文。
  \item 跑了 \texttt{/review} 就想一次性改完所有建议。
  \item 想自动化大范围改动，却还没在本地手工跑稳。
\end{itemize}

\subsection{MCP 先跑一个小例子}

\begin{lstlisting}[style=cli]
codex mcp list
codex mcp add docs --url https://example.com/mcp
\end{lstlisting}

最值得先试的，不是“全接上”，而是让 Codex 同时看本地代码和一个只读外部来源，比如产品文档或内部知识库。

\begin{tipbox}
第一次接 MCP，先做对照和解释，不要一上来就给外部写权限。
\end{tipbox}

\section{练习题}

\subsection{练习一}

让 Codex 给出一个公开小仓库的入口文件、核心目录和推荐阅读顺序。

\subsection{练习二}

选一个功能，例如登录或上传，让它解释入口、涉及文件和调用链。

\subsection{练习三}

只做一个低风险改动，例如改文案、补说明或补一条简单测试。

\subsection{练习四}

跑一遍：

\begin{lstlisting}[style=cli]
/review
/diff
\end{lstlisting}

然后判断它给出的审查结果，是否真的落到了当前改动上。

\subsection{练习五}

把一句模糊请求改写成更清楚的版本。比如把“帮我优化一下这里”改成包含目标、范围、限制、输出格式的提示词，再比较前后效果。

\subsection{做完后怎么判断合不合格？}

如果你做题前就想先看判分标准，也完全可以先读这一节。

\begin{itemize}
  \item 练习一至少要给出具体路径、推荐顺序、为什么先看这些文件。
  \item 练习二至少要讲清入口、关键文件和调用链，不只是列文件名。
  \item 练习三至少要说清改了哪些文件、没改什么、怎么验证。
  \item 练习四至少要把真正的问题和可选优化分开。
  \item 练习五至少要把目标、范围、限制、输出格式写清楚。
\end{itemize}

\section{最后只记这几句}

\begin{tcolorbox}[colback=SummaryLight,colframe=Summary,title=3 分钟背完,fonttitle=\bfseries]
\begin{enumerate}
  \item 日常开发多用 \texttt{codex}，脚本化多用 \texttt{codex exec}。
  \item 第一次上手先做最小任务，不要急着大改。
  \item 好提示词要写清楚目标、范围、限制、输出格式。
  \item 本地开发先从 \texttt{workspace-write} 起步。
  \item 改完一定先 \texttt{/review}，再 \texttt{/diff}，最后自己验证。
  \item 高风险任务先看计划、风险和回滚，再决定要不要放行。
  \item 有重复规则就写进 \texttt{AGENTS.md}。
  \item 查本地问题先看仓库，真要最新资料再开 \texttt{--search}。
\end{enumerate}
\end{tcolorbox}

\end{document}
